\documentclass[10pt]{article}

% COMP 2300 Counting Core companion: student-facing reference and practice.
\usepackage[letterpaper,margin=0.56in,includehead,includefoot,
  headheight=15pt,headsep=7pt,footskip=20pt]{geometry}
\usepackage[T1]{fontenc}
\usepackage{amsmath,amssymb,mathtools}
\usepackage{microtype}
\usepackage[dvipsnames,svgnames]{xcolor}
\usepackage{multicol}
\usepackage{enumitem}
\usepackage{tabularx,booktabs,array}
\usepackage[most]{tcolorbox}
\usepackage{fancyhdr}
\usepackage{titlesec}
\usepackage[hidelinks]{hyperref}

\definecolor{CoreNavy}{HTML}{18364A}
\definecolor{CoreTeal}{HTML}{176B6B}
\definecolor{CoreGold}{HTML}{9A6500}
\definecolor{CoreRed}{HTML}{963D38}
\definecolor{CoreBlueTint}{HTML}{EDF4F8}
\definecolor{CoreTealTint}{HTML}{EDF8F6}
\definecolor{CoreGoldTint}{HTML}{FFF7E5}
\definecolor{CoreRedTint}{HTML}{FCEFED}

\setlength{\parindent}{0pt}
\setlength{\parskip}{2.5pt}
\setlength{\columnsep}{0.25in}
\setlength{\columnseprule}{0.25pt}
\setlength{\emergencystretch}{1.2em}
\setlength{\abovedisplayskip}{3pt}
\setlength{\belowdisplayskip}{3pt}
\setlist{leftmargin=*,itemsep=1pt,topsep=1.5pt,parsep=0pt,partopsep=0pt}
\raggedbottom

\titleformat{\section}{\Large\bfseries\color{CoreNavy}}{}{0pt}{}
  [\vspace{-3pt}\color{CoreNavy}\rule{\linewidth}{0.7pt}]
\titlespacing*{\section}{0pt}{3pt}{4pt}

\pagestyle{fancy}
\fancyhf{}
\fancyhead[L]{\textbf{COMP 2300}}
\fancyhead[C]{\textbf{Counting Core Companion}}
\fancyhead[R]{Definitions and Worked Examples}
\fancyfoot[L]{Core sections A--E}
\fancyfoot[C]{\thepage}
\fancyfoot[R]{Study reference---not used during a Core}
\renewcommand{\headrulewidth}{0.5pt}
\renewcommand{\footrulewidth}{0.3pt}

\tcbset{enhanced,before skip=2.5pt,after skip=2.5pt,boxrule=0.55pt,arc=1.2mm,
  left=4pt,right=4pt,top=3pt,bottom=3pt,fonttitle=\bfseries}
\newtcolorbox{definitionbox}[1]{title={#1},colback=CoreBlueTint,
  colframe=CoreNavy,colbacktitle=CoreNavy,coltitle=white}
\newtcolorbox{methodbox}[1]{title={#1},colback=CoreTealTint,
  colframe=CoreTeal,colbacktitle=CoreTeal,coltitle=white}
\newtcolorbox{workedbox}[1]{title={#1},colback=white,
  colframe=CoreNavy,colbacktitle=CoreBlueTint,coltitle=CoreNavy}
\newtcolorbox{warningbox}[1]{title={#1},colback=CoreRedTint,
  colframe=CoreRed,colbacktitle=CoreRed,coltitle=white}
\newtcolorbox{checklistbox}[1]{title={#1},colback=CoreGoldTint,
  colframe=CoreGold,colbacktitle=CoreGold,coltitle=white}
\newcommand{\term}[1]{\textbf{\color{CoreNavy}#1}}

\begin{document}
\small

\section*{A--B. Counting Rules, Arrangements, and Selections}
\begin{multicols}{2}
\begin{definitionbox}{Sum and product rules}
Use the \term{sum rule} for disjoint alternatives: $m+n$ ways to do one of
two tasks. Use the \term{product rule} for a sequence of choices: $mn$ ways
when every first choice is followed by $n$ second choices. If alternatives
overlap, use inclusion--exclusion:
\[
|A\cup B|=|A|+|B|-|A\cap B|.
\]
\end{definitionbox}

\begin{workedbox}{A.1--A.4 worked examples}
\textbf{Disjoint alternatives:} five soups or four salads give $5+4=9$
lunches. \textbf{Successive positions:} two letters followed by three digits
give $26\cdot26\cdot10\cdot10\cdot10=26^2 10^3$ strings when repetition is
allowed. \textbf{Functions:} choose one of eight images for each of five
labeled inputs, giving $8^5$. \textbf{Overlap:} if $|A|=19$, $|B|=20$, and
$|A\cap B|=8$, add both sets and remove the duplicate intersection:
$|A\cup B|=19+20-8=31$.
\end{workedbox}

\begin{definitionbox}{Generalized pigeonhole principle}
Placing $N$ objects into $k$ boxes guarantees some box contains at least
\[
\left\lceil\frac Nk\right\rceil
\]
objects. Forty keys mapped to twenty hash values guarantee occupancy at least
$\lceil40/20\rceil=2$. They do not guarantee $3$: two keys per bucket is
possible.
\end{definitionbox}

\columnbreak

\begin{definitionbox}{Permutations and combinations}
When order matters, select and arrange:
\[
P(n,r)=n(n-1)\cdots(n-r+1)=\frac{n!}{(n-r)!}.
\]
When order does not matter, select:
\[
\binom nr=\frac{n!}{r!(n-r)!}.
\]
For repeated object types with multiplicities $n_1,\ldots,n_t$ totaling $n$,
the number of distinct arrangements is $n!/(n_1!\cdots n_t!)$.
\end{definitionbox}

\begin{workedbox}{B worked examples}
A committee is unordered, so three people from ten give $\binom{10}{3}$.
Offices are labeled, so president then vice president give
$P(10,2)=10\cdot9$. For \texttt{BANANA}, begin with $6!$ arrangements, then
remove the duplicate permutations of the three A's and two N's:
\[
\frac{6!}{3!2!}
\]
An eight-bit string with exactly three $1$s is determined by choosing the three
$1$-positions; the other five positions are forced, giving $\binom83$.
\end{workedbox}

\begin{definitionbox}{Binomial theorem}
\[
(x+y)^n=\sum_{k=0}^{n}\binom nk x^{n-k}y^k.
\]
The coefficient of $x^3$ in $(1+x)^8$ is $\binom83$: choose the three factors
that contribute $x$.
\end{definitionbox}

\begin{warningbox}{Frequent A--B errors}
Do not add sequential choices or multiply overlapping alternatives. Decide
whether changing order creates a new outcome before choosing $P(n,r)$ or
$\binom nr$. Divide only for genuinely indistinguishable repetitions.
\end{warningbox}

\begin{checklistbox}{A--B readiness}
Can I name the rule, show an unsimplified expression, explain order, and state
what each factor or selected position represents?
\end{checklistbox}
\end{multicols}

\clearpage
\section*{C. Constrained Counting}
\begin{multicols}{2}
\begin{methodbox}{Four reliable strategies}
\begin{enumerate}
  \item \term{Build directly}: multiply successive legal choices.
  \item \term{Split into disjoint cases}: count each case, then add.
  \item \term{Use a complement}: total minus the forbidden cases.
  \item \term{Choose positions or members}: use binomial coefficients, then
  fill remaining roles.
\end{enumerate}
State the sample objects and why the cases neither overlap nor leave gaps.
\end{methodbox}

\begin{workedbox}{C.1: exact composition}
\textbf{1. Translate ``exactly'':} a five-card hand with exactly two hearts has
two hearts and three nonhearts. \textbf{2. Choose within each category:}
\[
\binom{13}{2}\binom{39}{3}.
\]
\textbf{3. Justify the product:} each qualifying hand has one unique pair of
hearts and triple of nonhearts, so multiply the two unordered selections.
\end{workedbox}

\begin{workedbox}{C.2: count by complement}
\textbf{1. Count all:} there are $2^8$ eight-bit strings.
\textbf{2. Identify the complement:} ``at least two $1$s'' fails only for zero
or one $1$. \textbf{3. Subtract the disjoint forbidden cases:}
\[
2^8-\binom80-\binom81.
\]
The complement is shorter because it has only these two cases.
\end{workedbox}

\begin{workedbox}{C.3: mixed committee}
A five-person group uses disjoint pools of twelve CS majors and eight math
majors. Exactly three CS majors therefore requires two math majors. Choose
the two unordered subsets in
\[
\binom{12}{3}\binom82
\]
ways. ``Exactly'' fixes both category counts.
\end{workedbox}

\columnbreak

\begin{workedbox}{C.4: injective functions}
For $f:A\to B$ with $|A|=5$ and $|B|=8$, \textbf{1. label the five inputs.}
\textbf{2. Enforce injectivity:} after an output is used, it is unavailable.
The successive choices are
\[
8\cdot7\cdot6\cdot5\cdot4=P(8,5).
\]
\textbf{3. Check overcounting:} input labels already distinguish the positions,
so do not divide by $5!$.
\end{workedbox}

\begin{workedbox}{C.5: codes with a first-position rule}
\textbf{1. Enforce the special position first:} a six-letter code beginning
with a vowel has five first-character choices. \textbf{2. Remove that letter:}
without repetition, the remaining ordered choices are $25,24,23,22,21$:
\[
5\cdot25\cdot24\cdot23\cdot22\cdot21.
\]
Each factor now has a visible job: the first enforces the vowel rule, and the
others enforce no repetition.
\end{workedbox}

\begin{definitionbox}{Avoiding overcounting}
A counting method is valid only if each desired object is generated exactly
once. Ask:
\begin{itemize}
  \item Can the same outcome arise from two different choices or cases?
  \item Did I assign an order to an unordered object?
  \item Did I forget labels that make positions distinct?
  \item Are my cases exhaustive and disjoint?
\end{itemize}
\end{definitionbox}

\begin{warningbox}{Frequent C errors}
``At least'' usually needs several cases or a complement. ``Without
replacement'' changes later factors. A restriction on the first character
must be enforced before counting unrestricted remaining positions. Never use
$n^r$ for an injective function unless $r\le1$.
\end{warningbox}

\begin{checklistbox}{C readiness}
Can I construct the count two ways, identify possible overlap, justify every
factor, and leave a correct unsimplified expression?
\end{checklistbox}
\end{multicols}

\clearpage
\section*{D. Probability Models}
\begin{multicols}{2}
\begin{definitionbox}{Sample spaces and events}
A \term{sample space} $\Omega$ lists all possible outcomes; an \term{event}
is a subset of $\Omega$. For equally likely outcomes,
\[
P(A)=\frac{|A|}{|\Omega|}.
\]
Always justify equal likelihood. Useful identities are
\begin{align*}
P(A^c)&=1-P(A),\\
P(A\cup B)&=P(A)+P(B)-P(A\cap B),\\
P(A\cap B)&=P(A)P(B\mid A).
\end{align*}
\end{definitionbox}

\begin{workedbox}{D.1: coin and die}
\textbf{Model:} toss a fair coin and roll a fair die independently.
\textbf{Sample space:} there are $2\cdot6=12$ equally likely coin--die pairs.
\textbf{Favorable count:} heads gives $6$, an even die gives $6$, and their
intersection gives $3$. Hence $6+6-3=9$ favorable outcomes and
$P=9/12=3/4$. The subtraction removes the duplicate intersection.
\end{workedbox}

\begin{workedbox}{D.2: probability identities}
If $P(A)=0.55$, $P(B)=0.40$, and $P(A\cap B)=0.25$, then
\[
P(A\cup B)=0.55+0.40-0.25=0.70.
\]
Also $B=(A\cap B)\mathbin{\dot\cup}(A^c\cap B)$, so
$P(A^c\cap B)=0.40-0.25=0.15$.
\end{workedbox}

\begin{workedbox}{D.3: uniform bit strings}
\textbf{Denominator:} $2^8$ equally likely eight-bit strings.
\textbf{Numerator:} choose the three positions containing $1$s, so $\binom83$
strings qualify. Thus
\[
P(\text{exactly three }1\text{s})=\frac{\binom83}{2^8}.
\]
The numerator and denominator must describe the same sample space.
\end{workedbox}

\columnbreak

\begin{definitionbox}{Conditional probability and independence}
For $P(B)>0$,
\[
P(A\mid B)=\frac{P(A\cap B)}{P(B)}.
\]
$A$ and $B$ are independent exactly when
$P(A\cap B)=P(A)P(B)$, equivalently $P(A\mid B)=P(A)$ when defined.
Independence is a claim to test, not a default assumption.
\end{definitionbox}

\begin{workedbox}{D.4: without replacement}
A bag has four red and six blue tokens. \textbf{First draw:} red has
probability $4/10$. \textbf{Update:} after removing a red, nine tokens remain
and all six blue tokens remain. Therefore
\[
\frac4{10}\cdot\frac69.
\]
The second factor is conditional on the first draw; keeping denominator $10$
would describe replacement, not this experiment.
\end{workedbox}

\begin{workedbox}{D.5: two different conditionals}
If $24$ students use Python, $18$ use Java, and $10$ use both, then
\[
P(P\mid J)=\frac{10}{18},\qquad
P(J\mid P)=\frac{10}{24}.
\]
The shared numerator is the intersection. The event after the conditioning bar
determines the denominator, so reversing the conditional usually changes the
value.
\end{workedbox}

\begin{warningbox}{Frequent D errors}
Do not add probabilities for a non-disjoint ``or'' without subtracting the
overlap. Do not keep a denominator unchanged after sampling without
replacement. Do not confuse $P(A\mid B)$ with $P(B\mid A)$ or infer
independence from informal wording.
\end{warningbox}

\begin{checklistbox}{D readiness}
Can I name the sample space, choose the correct denominator, use a complement
or overlap identity, and distinguish conditional probability from
independence?
\end{checklistbox}
\end{multicols}

\clearpage
\section*{E. Modeling, Bayes, Expectation, and Interpretation}
\begin{multicols}{2}
\begin{definitionbox}{Total probability and Bayes' rule}
If $B_1,\ldots,B_m$ partition the sample space, then
\[
P(A)=\sum_i P(A\mid B_i)P(B_i).
\]
Bayes' rule reverses the condition:
\[
P(B_j\mid A)=\frac{P(A\mid B_j)P(B_j)}{P(A)}.
\]
The prior or base rate $P(B_j)$ remains part of the numerator.
\end{definitionbox}

\begin{workedbox}{E.1: numerical independence test}
If $P(A)=0.6$, $P(B)=0.5$, and $P(A\cap B)=0.3$, then
$P(A)P(B)=0.3=P(A\cap B)$, so the events are independent. Consequently
$P(A\mid B)=0.3/0.5=0.6=P(A)$.
\end{workedbox}

\begin{workedbox}{E.2: screening and base rates}
Let $C$ be the condition and $+$ a positive test. Then
\[
\underbrace{0.95(0.01)}_{\text{true positives}}=0.0095,
\qquad
\underbrace{0.08(0.99)}_{\text{false positives}}=0.0792.
\]
Add the two disjoint routes: $P(+)=0.0095+0.0792=0.0887$. Then Bayes' rule
keeps the true-positive route in the numerator:
\[
P(C\mid +)=\frac{0.0095}{0.0887}\approx0.107.
\]
Even a sensitive test can yield many false positives when the condition is
rare; the denominator includes both true and false positives.
\end{workedbox}

\begin{workedbox}{E.3: factory Bayes model}
Factory A supplies $40\%$ with defect rate $3\%$; B supplies $60\%$ with rate
$1\%$. First compute both joint routes:
\[
\begin{aligned}
P(A\cap D)&=0.40(0.03)=0.012,\\
P(B\cap D)&=0.60(0.01)=0.006.
\end{aligned}
\]
Thus $P(D)=0.018$ and $P(A\mid D)=0.012/0.018=2/3$.
\end{workedbox}

\columnbreak

\begin{definitionbox}{Expected value}
For numerical outcomes $x_i$ with probabilities $p_i$,
\[
E[X]=\sum_i x_i p_i.
\]
Use \emph{net} outcomes when a cost is paid. Expected value is a long-run
average, not a guaranteed result from one trial.
\end{definitionbox}

\begin{workedbox}{E.4: expected net value}
A fair die pays $10$ on a $6$, $2$ on a $4$ or $5$, and $0$ otherwise. With a
$2$ cost, first convert payouts to net outcomes: $8,0,-2$, with probabilities
$1/6,2/6,3/6$. Therefore
\[
E[X]=8\left(\frac16\right)+0\left(\frac26\right)
-2\left(\frac36\right)=\frac13.
\]
The positive sign means an average gain of $\$1/3$ per play over many plays.
\end{workedbox}

\begin{definitionbox}{Aggregation reversal}
A treatment can have the higher success rate in every subgroup but a lower
overall rate when the groups receive treatments in very different proportions.
Compute each within-group rate and each pooled rate separately. Then explain
which high- or low-success group dominates each pooled denominator; do not say
that arithmetic is contradictory.
\end{definitionbox}

\begin{methodbox}{Twenty-minute model-selection routine}
For one problem from each A--E section: name the rule or model first, write an
exact unsimplified expression, then calculate or interpret. Check that each
denominator matches its condition and that each counted outcome appears once.
\end{methodbox}

\begin{warningbox}{Frequent E errors}
Do not ignore base rates, reverse a conditional, or assume independence. In
expected value, subtract costs from every applicable outcome. In aggregated
tables, compare rates rather than raw success counts.
\end{warningbox}

\begin{checklistbox}{Whole-Core readiness}
I can choose among sum, product, cases, complement, permutation, and
combination; justify a probability denominator; compute conditional and Bayes
probabilities; test independence; calculate net expected value; and explain an
aggregation reversal. A live form contains one problem from each A--E section.
\end{checklistbox}
\end{multicols}

\end{document}
